% Independently audited extension, 7 October 2026.
% Uses manuscript macros \ph, \stp, \detp, \rd, and theorem environment.
% Add bibliography key dritschelwoerdeman as recorded below.
\section{The sharp logarithmic degree bound for four equal masses}
\label{sec:logupper}
The linear upper bound can be sharpened by using both competing hierarchy choices. The resulting bound is uniform even when the posterior gaps degenerate.

\begin{theorem}[Sharp degree order for four equal masses]
\label{thm:logupper}
There is an absolute constant $C$ such that, for every $r\ge1$, every nonnegative nonzero polynomial curvature $g=\ph''$ of degree at most $r$, and every four distinct equally weighted posterior locations,
\[
\rho_{\ph,\stp}\le\rho_{\ph,\detp}
\le C\frac{r+1}{\log(r+2)}.
\]
Consequently, together with the fixed-quartet lower construction in Theorem~\ref{thm:lower},
\[
C_r^{\detp}=\Theta(r/\log r),\qquad
C_r^{\stp}=\Theta(r/\log r)\qquad(r\to\infty).
\]
The upper bound allows arbitrary posterior spacing; the lower bound already holds for one fixed equally spaced quartet. This is an asymptotic-order result, with no claim of optimal leading constants or historical novelty.
\end{theorem}
\begin{proof}
Write $R=\rho_{\ph,\detp}$. The compatible-optimizer argument of Theorem~\ref{thm:linearupper} shows that either $R=1$, or the only relevant conflict is the flat fine pair $BC$ against the flat coarse split $AB\mid CD$, with $A<B<C<D$. Put
\[
q=P_{BC},\qquad Q=P_{AB}+P_{CD},\qquad T=T_{ABC}.
\]
The hierarchy keeping the optimal coarse split and its better outer fine pair has price $\min(P_{AB},P_{CD})/q$. The hierarchy keeping the optimal fine pair and its better adjacent triple has price $\min(T_{ABC},T_{BCD})/Q$. Therefore
\[
R\le\frac{\min(P_{AB},P_{CD})}{q},\qquad
R\le\frac{\min(T_{ABC},T_{BCD})}{Q}.
\]
These are merely admissible candidates: no exact Pareto-frontier assertion, or exclusion of a noncontiguous improvement, is needed. In particular,
\begin{equation}
q\le \frac{T}{R^2},\qquad
S:=P_{AB}+P_{BC}\le\frac{2T}{R}.
\label{eq:logbothratios}
\end{equation}

Affinely normalize the \emph{triple hull} $[A,C]$ to $[-1,1]$, writing $B=b\in(-1,1)$. The fourth point need not remain in this interval; it is not used in the polynomial estimates below. Set
\[
w(t)=1-|t|,\quad w_b=w(b),\quad
s=\sqrt{1-b^2},\quad \theta_b=\arccos b,\quad
\mu=\int_{-1}^1w(t)g(t)\rd t.
\]
The triple kernel dominates the endpoint-pair kernel $w/4$, and is at most $w/2$ because it vanishes at both endpoints and has piecewise slope of magnitude at most $1/2$. Hence
\begin{equation}
\frac{w(t)}4\le K_{ABC}(t)\le\frac{w(t)}2,
\qquad 2T\le\mu\le4T.
\label{eq:logtriplekernel}
\end{equation}
The two adjacent-pair kernels sum to
\[
\frac14\min\{w(t),|t-b|\}.
\]

\emph{A peak near the inner knot.}
Assume $R\ge4096$, and put
\[
\varepsilon=\frac{64}{R}\le\frac1{64},\qquad
I=\{t:|t-b|\le\varepsilon w_b\}.
\]
Outside $I$, the 1-Lipschitz property of $w$ gives
$|t-b|\ge\varepsilon w(t)/(1+\varepsilon)$. Thus
\[
\int_{[-1,1]\setminus I}wg
\le\frac{4(1+\varepsilon)}{\varepsilon}S
\le\frac{8(1+\varepsilon)T}{R\varepsilon}
\le\frac T4.
\]
It follows that $\int_Iwg\ge T$, so some $t_0\in I$ satisfies
$w(t_0)g(t_0)\ge T/(2\varepsilon w_b)$. Define
\[
F(\theta)=g(\cos\theta)\sin^4\theta,\qquad
n=r+4,\qquad \theta_0=\arccos t_0.
\]
This is a nonnegative trigonometric polynomial of degree at most $n$. Since $(1-t^2)^2\ge w(t)^2$,
\begin{equation}
F(\theta_0)\ge\frac{T}{4\varepsilon}.
\label{eq:logpeak}
\end{equation}
Put $h=\varepsilon s$ and $a=1-h$. Along the interval between $b$ and $t_0$,
$|(1-t^2)-s^2|\le2\varepsilon w_b\le2\varepsilon s^2$. Consequently,
\begin{equation}
|\theta_0-\theta_b|\le2h,\qquad \sin\theta_0\le2s.
\label{eq:logangles}
\end{equation}
Let $\nu$ be the Poisson probability measure centered at $ae^{i\theta_0}$:
\[
\rd\nu(\theta)=P_a(\theta-\theta_0)\frac{\rd\theta}{2\pi},
\qquad
P_a(u)=\frac{1-a^2}{1-2a\cos u+a^2}.
\]

\emph{A spacing-uniform Poisson average.}
For circular distance $d=\operatorname{dist}(u,2\pi\mathbb Z)\in[0,\pi]$, the facts $a\ge63/64$ and $1-\cos u\ge2d^2/\pi^2$ imply
\[
P_a(u)\le\frac{8h}{h^2+d^2}.
\]
For $\theta\in[0,\pi]$, either circular distance from $\theta$ to $\theta_0$ or to $-\theta_0$ satisfies
$\sin\theta\le\sin\theta_0+d\le2s+d$. Therefore
\[
P_a(\theta\pm\theta_0)\sin\theta
\le\frac{16hs}{h^2+d^2}+\frac{8hd}{h^2+d^2}
\le\frac{20}{\varepsilon}.
\]
Using evenness of $F$, changing variables $t=\cos\theta$, and noting
$(1-t^2)^{3/2}/w(t)\le2\sin\theta$, we obtain
\begin{align}
\int F\rd\nu
&=\frac1{2\pi}\int_{-1}^1
 [P_a(\theta-\theta_0)+P_a(\theta+\theta_0)]
 g(t)(1-t^2)^{3/2}\rd t \notag\\
&\le\frac{80\mu}{2\pi\varepsilon}
\le\frac{64T}{\varepsilon}.
\label{eq:logpoissonwhole}
\end{align}
Here $\theta=\arccos t$. This estimate avoids any need to compare the unweighted angular mass of $F$ to the mass near $b$ when $b$ approaches an endpoint.

\emph{An inward arc with little curvature.}
Consider
\[
J=[\theta_b-5h,\theta_b-4h].
\]
It lies in $(0,\pi)$, since $s\le\theta_b$ and $5\varepsilon<1$. The sine function is 1-Lipschitz, so on $J$, and on the interval joining it to $\theta_b$,
\[
(1-5\varepsilon)s\le\sin\theta\le(1+5\varepsilon)s.
\]
For $t=\cos\theta$ on this arc, this gives
\[
3\varepsilon w_b\le3\varepsilon s^2\le t-b
\le6\varepsilon s^2\le12\varepsilon w_b,
\qquad
1-t\ge(1-12\varepsilon)w_b\ge3\varepsilon w_b.
\]
The $BC$ pair kernel is therefore at least $3\varepsilon w_b/4$. Using $s^2/w_b\le2$,
\[
\int_JF(\theta)\rd\theta
\le8s^3\int_{\cos J}g(t)\rd t
\le\frac{64s}{3\varepsilon}q.
\]
Since $P_a\le2/h$, \eqref{eq:logbothratios} yields
\begin{equation}
\int_JF\rd\nu
\le\frac{64q}{3\pi\varepsilon^2}
\le\frac{T}{192\pi}\le T.
\label{eq:logleakage}
\end{equation}
Write $\alpha=\nu(J)$. By \eqref{eq:logangles},
$|\theta-\theta_0|\le7h$ on $J$. The denominator of $P_a$ is then at most $50h^2$, whereas its numerator is at least $h$. Because $|J|=h$,
\begin{equation}
\alpha\ge\frac1{100\pi}=:\alpha_0.
\label{eq:logharmonicmass}
\end{equation}

\emph{Outer factorization forces the logarithm.}
The classical Fej\'er--Riesz theorem supplies an algebraic polynomial $p$ of degree at most $n$, with no zeros in the open unit disk, such that
$F(\theta)=|p(e^{i\theta})|^2$; see \cite[Introduction]{dritschelwoerdeman}. Unit-circle zeros are allowed. For $|z|=1$ and any root $\zeta$ with $|\zeta|\ge1$,
\[
|az-\zeta|^2-a^2|z-\zeta|^2
=(1-a)\big[(1+a)|\zeta|^2-2a\operatorname{Re}(z\overline\zeta)\big]\ge0.
\]
Multiplication over the roots proves
$|p(ae^{i\theta_0})|\ge a^n|p(e^{i\theta_0})|$.
The Poisson formula for $\log|p|$ therefore gives
\begin{equation}
\log F(\theta_0)+2n\log a\le\int\log F\rd\nu.
\label{eq:logouter}
\end{equation}
For completeness, the formula can be checked factor by factor. For a root strictly outside the disk, $\log|z-\zeta|$ is harmonic on a neighborhood of the closed disk. For a unit-circle root, replace $\zeta$ by $(1+\eta)\zeta$ and let $\eta\downarrow0$. The boundary logarithms converge in $L^1$: locally they are bounded below by the integrable logarithm of the unit-circle distance and above by a fixed constant. The Poisson kernel at the fixed interior point is bounded, so its integrals converge as well. Thus boundary zeros do not add any correction to the formula.

Apply Jensen's inequality separately on $J$ and its complement. From
\eqref{eq:logpoissonwhole}--\eqref{eq:logharmonicmass} and the binary entropy bound,
\begin{align*}
\int\log F\rd\nu
&\le\alpha\log\frac{T}{\alpha}
 +(1-\alpha)\log\frac{64T}{(1-\alpha)\varepsilon}\\
&\le\log\frac{T}{\varepsilon}+\alpha_0\log\varepsilon+\log128.
\end{align*}
Combining with \eqref{eq:logpeak} and \eqref{eq:logouter}, then using
$-\log(1-h)\le2h\le2\varepsilon$, yields the explicit inequality
\begin{equation}
4(r+4)\varepsilon
\ge\frac1{100\pi}\log\frac1\varepsilon-\log512,
\qquad \varepsilon=\frac{64}{R},\quad R\ge4096.
\label{eq:logimplicit}
\end{equation}
Equivalently, with absolute constants
$K=64\cdot512^{100\pi}$ and $A=25600\pi$,
\[
R\log(R/K)\le A(r+4).
\]
For $R\ge\max(4096,K^2)$ this gives $R\log R\le2A(r+4)$ and hence
$R=O(r/\log r)$; the remaining fixed thresholds are absorbed into a universal constant. More explicitly, split into $R\le\sqrt{r+4}$ and its complement, on which $\log R\ge\tfrac12\log(r+4)$. The stochastic upper bound follows from deterministic admissibility, and the Chebyshev family in Theorem~\ref{thm:lower} supplies both matching lower orders.
\end{proof}

